\documentclass[../../../../main.tex]{subfiles}
\begin{document}

\begin{flushright}
\footnotesize\textit{【自動文字起こし・要確認】}
\end{flushright}

[解] $A$から北へ$n$メートルいった
点を$A_n$、その1m西を$B_n$とする

\begin{center}
\begin{tikzpicture}
  \draw[->] (-2,0) -- (2,0) node[right] {$x$};
  \draw (0,-0.2) -- (0,3);
  \fill (0,0) circle (1.5pt) node[below right] {$A$};
  \fill (-1,0) circle (1.5pt) node[below left] {$B$};
  \fill (0,2) circle (1.5pt) node[right] {$A_n$};
  \fill (-1,2) circle (1.5pt) node[left] {$B_n$};
  \draw[dashed] (-1,0) -- (-1,2) -- (0,2);
  \foreach \y in {0,1,2} {
    \draw (-0.1,\y) -- (0.1,\y);
    \draw (-1.1,\y) -- (-0.9,\y);
  }
\end{tikzpicture}
\end{center}

(1) $A_n$にたどりつくのは、$B_n$について
オモテ $\dots$ 右へ$1$
ウラ $\dots$ 上へ$1$
から表をだす確率で、
$$ p_n = \frac{{}_{n+4}C_4}{2^{n+4}} \cdot \frac{1}{2} = \frac{{}_{n+4}C_4}{2^{n+5}} $$

(2) $\frac{p_{n+1}}{p_n} \ge 1 \cdots (1)$ をみたす$n$をもとめる。(1)から、
$$ \frac{p_{n+1}}{p_n} = \frac{{}_{n+5}C_4}{{}_{n+4}C_4} \cdot \frac{2^{n+5}}{2^{n+6}} = \frac{(n+5)(n+4)(n+3)(n+2)}{(n+4)(n+3)(n+2)(n+1)} \frac{1}{2} $$
$$ = \frac{n+5}{2(n+1)} \ge 1 $$
$$ \iff 3 \ge n \quad (n=3\text{で等号成立}) $$
だから、
$$ p_1 < p_2 < p_3 = p_4 > p_5 > \cdots $$
より、$p_n$を最大にする$n$は、$n=3, 4$

\end{document}
